\section{Method}

\subsection{Multi-Route Variable-Depth Retrieval}
\label{sec:retrieval-method}

Figure~\ref{fig:retrieval} shows our Multi-Route Variable-Depth
Retrieval with Deep-Tail Reranking system. The original narrative
remains the primary retrieval signal, with separate budgets for
expansion, reranking, and serialization. Candidate membership,
ordering, and submitted depth are computed in separate stages.

\begin{figure*}[t]
\centering
\resizebox{0.95\textwidth}{!}{% Publication-oriented schematic of the final Retrieval pipeline. Query2Doc
% creates one expanded BM25 query (one pseudo-document plus the original
% narrative repeated five times); follow-up routes are added by the bounded
% acquisition loop. The cutoff is computed before deep-tail fusion and applied
% to the final post-deep ranking.
\definecolor{cfdaNavy}{HTML}{18324B}
\definecolor{cfdaSlate}{HTML}{526675}
\definecolor{cfdaBlueFill}{HTML}{EEF3F7}
\definecolor{cfdaBlueLine}{HTML}{607F99}
\definecolor{cfdaSageFill}{HTML}{EDF6F1}
\definecolor{cfdaSageLine}{HTML}{4F7F6F}
\definecolor{cfdaSandFill}{HTML}{FAF5E8}
\definecolor{cfdaSandLine}{HTML}{9A7A3D}
\definecolor{cfdaNeutral}{HTML}{F7F8F8}
\definecolor{cfdaFrame}{HTML}{CCD8E1}

\begin{tikzpicture}[
  box/.style={draw=cfdaSlate, rounded corners=2pt, fill=white,
    align=center, minimum height=0.88cm, inner sep=3pt, font=\small},
  process/.style={box, fill=cfdaBlueFill, draw=cfdaBlueLine},
  endpoint/.style={box, fill=cfdaSageFill, draw=cfdaSageLine,
    rounded corners=5pt},
  artifact/.style={box, fill=cfdaSandFill, draw=cfdaSandLine,
    rounded corners=4pt},
  input/.style={box, fill=cfdaNeutral},
  decision/.style={draw=cfdaBlueLine, diamond, aspect=2.05,
    fill=cfdaBlueFill, align=center, inner sep=1.8pt, font=\small},
  stage/.style={font=\small\bfseries, text=cfdaNavy},
  flow/.style={-{Stealth[length=2mm]}, line width=0.62pt, draw=cfdaSlate},
  loop/.style={flow, draw=cfdaBlueLine},
  signal/.style={-{Stealth[length=1.8mm]}, line width=0.56pt,
    draw=cfdaBlueLine, dashed},
  note/.style={font=\scriptsize, text=cfdaSlate, fill=white, inner sep=1pt},
  detail/.style={draw=cfdaFrame, rounded corners=2pt, fill=cfdaNeutral,
    align=center, font=\footnotesize, inner sep=3pt, text=cfdaSlate},
]

% Stage 1: initial routes and bounded follow-up acquisition.
\node[stage,anchor=west] (s1title) at (0.10,7.05)
  {Stage 1: agentic multi-route candidate generation};
\node[endpoint,text width=1.35cm] (narr) at (0.75,5.35) {Narrative};
\node[process,text width=4.35cm,minimum height=1.65cm] (initialroutes) at (4.25,5.35)
  {Initial BM25 routes\\Original narrative: weight 1.0\\Query2Doc: 1 pseudo-doc + narrative $\times5$\\weight 1.0; omit if unusable};
\node[process,text width=1.82cm] (rrf) at (7.95,5.35)
  {Initial weighted RRF\\$k=60$};
\node[box,text width=1.65cm] (read) at (10.30,5.35)
  {Initial read\\first 5 successfully\\fetched};
\node[decision] (judge) at (13.15,5.35) {Evidence\\sufficient?};
\node[endpoint,text width=2.02cm] (pool) at (18.45,5.35)
  {Final pre-rerank pool\\up to 5,000};

\draw[flow] (narr) -- (initialroutes);
\draw[flow] (initialroutes) -- (rrf);
\draw[flow] (rrf) -- (read);
\draw[flow] (read) -- (judge);
\draw[flow] (judge) -- (pool)
  node[midway,above,yshift=1.5pt,note] {sufficient / budget stop};

\node[box,text width=1.95cm] (validate) at (13.15,3.20)
  {Accept up to 3 valid\\follow-up queries};
\node[process,text width=2.15cm] (followup) at (16.05,3.20)
  {Separate BM25 routes\\depth 5,000\\weight 0.25 each};
\node[process,text width=1.82cm] (refuse) at (18.75,3.20)
  {Recompute\\weighted RRF};
\node[box,text width=1.62cm] (readmore) at (21.10,3.20)
  {Read up to\\$+4$ unseen};
\draw[loop] (judge.south) -- (validate.north)
  node[midway,right,note] {insufficient};
\draw[loop] (validate) -- (followup);
\draw[loop] (followup) -- (refuse);
\draw[loop] (refuse) -- (readmore);
\coordinate (loopright) at (22.35,3.20);
\coordinate (looptop) at (22.35,6.45);
\coordinate (loopin) at (13.15,6.45);
\draw[loop] (readmore.east) -- (loopright) -- (looptop) -- (loopin) -- (judge.north);
\node[note] at (17.75,6.45) {next acquisition iteration};
\node[detail,text width=6.40cm] (budgetnote) at (5.20,3.20)
  {cap 12 successfully read documents /\\max 3 acquisition iterations};
\node[inner sep=0pt] (loopbound) at (22.35,6.45) {};

% Stage 2: final ranking construction and serialization.
\node[stage,anchor=west] (s2title) at (0.10,1.15)
  {Stage 2: ranking and variable-depth serialization};
\node[input,text width=1.55cm] (poolin) at (0.85,-1.20) {Pre-rerank\\pool};
\node[process,text width=2.45cm,minimum height=1.55cm] (head) at (3.55,-1.20)
  {Top-100 head fusion\\BM25 + CE + dense\\RRF $1{:}1{:}1$, $k=60$};
\node[process,text width=2.55cm,minimum height=1.55cm] (facet) at (6.85,-1.20)
  {Facet retrieval\\$\leq20$ facets / $\leq40$ queries\\BM25@1,000; RRF $k=60$};
\node[process,text width=2.20cm,minimum height=1.55cm] (merge) at (9.95,-1.20)
  {Protected splice\\keep exact pre-facet\\Top-200};
\node[process,text width=2.60cm,minimum height=1.55cm] (deep) at (13.20,-1.20)
  {Deep-tail CE/pool fusion\\RRF $1{:}1$, $k=60$\\Dev 101--5,000\\Official 101--3,000};
\node[process,text width=1.85cm,minimum height=1.55cm] (select) at (16.50,-1.20)
  {Take top $k_q$\\of final ranking};
\node[endpoint,text width=1.52cm] (out) at (19.10,-1.20)
  {Six-column\\TREC run};
\node[detail,text width=5.30cm] (cutoff) at (13.25,0.72)
  {Compute $k_q$ from post-splice/pre-deep score\\$\tau=0.20$, rank-1 anchor; clip $[4,K]$\\Dev $K=1{,}000$ / Official $K=5{,}000$};
\node[detail,text width=9.20cm] (facetnote) at (5.65,-3.02)
  {tail: score $\downarrow$, best rank $\uparrow$, document ID lexicographic;\\remove Top-200 IDs; append unused pre-facet tail; failure $\rightarrow$ pre-facet ranking};
\node[detail,text width=7.30cm] (deepnote) at (14.65,-3.02)
  {Top-100 unchanged; official 3,001--5,000 retain pool order;\\membership unchanged};

\draw[flow] (poolin) -- (head);
\draw[flow] (head) -- (facet);
\draw[flow] (facet) -- (merge);
\draw[flow] (merge) -- (deep);
\draw[flow] (deep) -- (select);
\draw[flow] (select) -- (out);
\draw[signal] (merge.north) |- (cutoff.west);
\coordinate (depthturn) at (16.50,0.72);
\draw[signal,-] (cutoff.east) -- (depthturn);
\draw[signal] (depthturn) -- (select.north)
  node[midway,right=2.5pt,note] {depth $k_q$};

\begin{scope}[on background layer]
  \node[draw=cfdaFrame, fill=cfdaBlueFill!28, rounded corners=3pt,
    fit=(s1title)(narr)(initialroutes)(rrf)(read)(judge)(pool)
      (validate)(followup)(refuse)(readmore)(budgetnote)(loopbound), inner sep=8pt] {};
  \node[draw=cfdaFrame, fill=cfdaBlueFill!28, rounded corners=3pt,
    fit=(s2title)(poolin)(head)(facet)(merge)(deep)(select)(out)(cutoff)
      (facetnote)(deepnote), inner sep=8pt] {};
\end{scope}
\end{tikzpicture}
}
\caption{Multi-Route Variable-Depth Retrieval with Deep-Tail Reranking.
The diagram depicts the development configuration and separates
multi-route acquisition, head reranking, facet-based expansion,
deep-tail ordering, and variable-depth serialization. For the official
119-topic run, the deep-tail span was 101--3,000 and the serialization
cap was $K=5{,}000$, as specified in the text.}
\label{fig:retrieval}
\end{figure*}

The system begins with the original narrative. A Query2Doc route
generates one pseudo-document and appends it to an anchor-boosted query
formed by repeating the original narrative five times, expanding the
query without discarding the anchor
\citep{wang2023query2doc}. The original-narrative and successful
Query2Doc BM25 rankings are fused first (weight 1.0 each; RRF $k=60$),
and the loop reads the first five successfully fetched documents. Only
an insufficient judgment elicits follow-up queries. Each accepted query
retrieves BM25 to depth 5,000 as a new route of weight 0.25; fusion is
recomputed and up to four unseen documents are read. The loop stops on
sufficiency, at 12 documents, or after three acquisition iterations;
the fused pool contains up to 5,000 documents.

Retained Query2Doc traces use one temperature-zero call, a 400-token
limit, and a 180-word pseudo-document cap. A failed or degenerate call
omits that route; retrieval continues with the original narrative. Run configurations
select \texttt{gpt-5.6-sol} for Query2Doc and facets. The
\texttt{gpt-oss-120b} sufficiency judge instead proposes follow-ups,
with up to five attempts per decision. Topic records retain generated
Query2Doc and facet content, but no immutable provider snapshot or
per-response model identifier.

For a document $d$, weighted reciprocal rank fusion is
\begin{equation}
s_{\mathrm{RRF}}(d) = \sum_{r:\,d\in r} \frac{w_r}{k + \operatorname{rank}_r(d)}.
\label{eq:wrrf}
\end{equation}
Here $k=60$, the original and Query2Doc routes have weight 1.0, and
each accepted follow-up BM25 route has weight 0.25.
Equation~\ref{eq:wrrf} applies to
the Retrieval route-fusion stage; the separate top-300 RAG reranker is
defined in \S\ref{sec:rag-method}.

The fused candidate pool is built to depth 5,000. The top-100 head is
then reranked using BM25-rank, MiniLM cross-encoder-rank
\citep{wang2020minilm}, and BGE-M3 dense-embedding-rank \citep{chen2024bgem3} signals with
equal RRF weights and $k=60$. The top portion of the pre-facet ranking
is protected while facet-based retrieval constructs the complementary
tail. Up to 20 schema-validated, searchable facets tied to exact
narrative spans each yield one initial query and up to one reformulation
that does not use retrieved evidence. After normalized deduplication, each of at most
40 queries retrieves BM25 to depth 1,000. Equal-weight RRF ($k=60$)
accumulates repeated document IDs into a unique pool of up to 5,000.
The exact pre-facet top-200 remains unchanged; after its IDs are removed,
facet candidates are sorted by descending RRF score, ascending best
contributing rank, and then lexicographic document ID. They fill the tail, followed when
needed by unused pre-facet tail documents in original order. Failed routes are omitted; if
generation or every retrieval fails, the pre-facet ranking is used.

\paragraph{Deep cross-encoder tail reranking.}

The selected deep-tail stage applies
\url{cross-encoder/ms-marco-MiniLM-L6-v2} below the first 100
entries, which are preserved unchanged. It fuses cross-encoder and
candidate-pool tail ranks with equal RRF weights and $k=60$. The
reranked span is ranks 101--5,000 in development; in the official
119-topic run it is ranks 101--3,000, after which ranks 3,001--5,000
retain candidate-pool order. This stage does not add or remove
candidates, but it can change their tail order and which documents fall
inside a serialized prefix.

\paragraph{Variable-depth serialization.}

The topic-specific submission depth is derived from the post-splice,
pre-deep-CE candidate-pool score: the implementation specifies a relative threshold of
0.20, a minimum depth of 4, and a maximum depth $K$ that differs
between settings: $K=1{,}000$ for the development runs and $K=5{,}000$
for both official 119-topic submissions. The lower clip is never
active in either setting; the smallest depth selected by the threshold
is 100 on the development set and 137 on the official topics. Any such
cap is a team-imposed limit, as the Retrieval task specification sets
no fixed maximum submission depth. The deep-CE
stage changes ordering but does not change topic depth. The output is
serialized as a six-column TREC run. For topic $q$, the corresponding
serialization rule is
\begin{equation}
\begin{aligned}
k_q &= \operatorname{clip}\bigl(
\lvert\{d : s_q(d) \geq \tau\, s_q(d_q^{(1)})\}\rvert,\\[-2pt]
&\hspace{25mm}4, K\bigr),
\end{aligned}
\label{eq:vark}
\end{equation}
where $s_q(d)$ is the post-splice, pre-deep-CE candidate-pool score,
$d_q^{(1)}$ is the document at rank 1 of topic $q$, $\tau=0.20$, and
$K$ is the maximum-depth cap defined above.
The threshold determines the topic-specific submission depth $k_q$;
after deep-tail reranking, the system submits the top-$k_q$ prefix of
the final ranking.
The anchor is the rank-1 score rather than the pool maximum. Because
the spliced ranking can contain entries whose scores originate from
different fusion computations, $s_q$ is not necessarily monotone in
rank, and the maximum-scoring document can occur below the protected
splice boundary. Using the
maximum instead of rank 1 changes $k_q$ for 88 of the 119 official
topics. The cutoff is computed before deep-tail reranking. The deep
stage leaves $k_q$ unchanged but can alter both the order and
membership of the submitted prefix.
